\begin{exercises}
  \recommended \item
    求每个向量的规范名称。
    \begin{exparts}
      \partsitem \( \Re^2 \) 中从 \( (2,1) \) 到 \( (4,2) \) 的向量
      \partsitem \( \Re^2 \) 中从 \( (3,3) \) 到 \( (2,5) \) 的向量
      \partsitem \( \Re^3 \) 中从 \( (1,0,6) \) 到 \( (5,0,3) \) 的向量
      \partsitem \( \Re^3 \) 中从 \( (6,8,8) \) 到 \( (6,8,8) \) 的向量
    \end{exparts}
    \begin{answer}
      \begin{exparts*}
        \partsitem \( \colvec[r]{2 \\ 1}  \)
        \partsitem \( \colvec[r]{-1 \\ 2}  \)
        \partsitem \( \colvec[r]{4 \\ 0 \\ -3}  \)
        \partsitem \( \colvec[r]{0 \\ 0 \\ 0}  \)
      \end{exparts*}  
     \end{answer}
  \recommended \item 
    判断这两个向量是否相等。
    \begin{exparts}
      \partsitem 从 \( (5,3) \) 到 \( (6,2) \) 的向量与
        从 \( (1,-2) \) 到 \( (1,1) \) 的向量
      \partsitem 从 \( (2,1,1) \) 到 \( (3,0,4) \) 的向量与
        从 \( (5,1,4) \) 到 \( (6,0,7) \) 的向量
    \end{exparts}
    \begin{answer}
      \begin{exparts}
        \partsitem 不相等，它们的规范位置不同。
          \begin{equation*}
            \colvec[r]{1 \\ -1}
            \qquad
            \colvec[r]{0 \\ 3}
          \end{equation*}
        \partsitem 相等，它们的规范位置相同。
          \begin{equation*}
            \colvec[r]{1 \\ -1 \\ 3}
          \end{equation*}
      \end{exparts}  
     \end{answer}
  \recommended \item
    \( (1,0,2,1) \) 是否在经过
    \( (-2,1,1,0) \) 和 \( (5,10,-1,4) \) 的直线上？
    \begin{answer}
      这条直线是这个集合。
      \begin{equation*}
        \set{\colvec[r]{-2 \\ 1 \\ 1 \\ 0}
             +\colvec[r]{7 \\ 9 \\ -2 \\ 4}t \suchthat t\in\Re }
      \end{equation*}
      注意，方程组
      \begin{equation*}
        \begin{linsys}{2}
          -2  &+  &7t  &=  &1  \\
           1  &+  &9t  &=  &0  \\
           1  &-  &2t  &=  &2  \\
           0  &+  &4t  &=  &1  
        \end{linsys}
      \end{equation*}
      无解。
      因此所给的点不在这条直线上。  
    \end{answer}
  \recommended \item
     \begin{exparts}
        \partsitem 描述经过 \( (1,1,5,-1) \)、
           \( (2,2,2,0) \) 与 \( (3,1,0,4) \) 的平面。
        \partsitem 原点在那个平面内吗？
      \end{exparts}
    \begin{answer}
      \begin{exparts}
        \partsitem 注意到
          \begin{equation*}
            \colvec[r]{2 \\ 2 \\ 2 \\ 0}
            -\colvec[r]{1 \\ 1 \\ 5 \\ -1}
            =\colvec[r]{1 \\ 1 \\ -3 \\ 1}
            \qquad
            \colvec[r]{3 \\ 1 \\ 0 \\ 4}
            -\colvec[r]{1 \\ 1 \\ 5 \\ -1}
            =\colvec[r]{2 \\ 0 \\ -5 \\ 5}
          \end{equation*}
          所以该平面是这个集合。
          \begin{equation*}
            \set{\colvec[r]{1 \\ 1 \\ 5 \\ -1}
                 +\colvec[r]{1 \\ 1 \\ -3 \\ 1}t
                 +\colvec[r]{2 \\ 0 \\ -5 \\ 5}s
                \suchthat t,s\in\Re}
          \end{equation*}
        \partsitem 不在；方程组
          \begin{equation*}
            \begin{linsys}{3}
              1  &+  &1t  &+  &2s  &=  &0  \\
              1  &+  &1t  &   &    &=  &0  \\
              5  &-  &3t  &-  &5s  &=  &0  \\
             -1  &+  &1t  &+  &5s  &=  &0  
            \end{linsys}
          \end{equation*}
          无解。
      \end{exparts}  
    \end{answer}
  \item 给出下面每一个的向量描述。
    \begin{exparts}
      \item $\Re^3$~中方程为 $x-2y+z=4$ 的平面子集
      \item $\Re^3$~中方程为 $2x+y+4z=-1$ 的平面
      \item $\Re^4$~中方程为 $x+y+z+w=10$ 的超平面子集
    \end{exparts}
    \begin{answer}
      \begin{exparts}
        \item 把 $x-2y+z=4$ 看作只含一个方程的线性方程组，
          以变量 $y$ 和~$z$ 为参数进行参数化，得到
          $x=4+2y-z$。
          由此得到这个平面的向量描述。
          \begin{equation*}
            \set{\colvec{x \\ y \\ z}=\colvec{4 \\ 0 \\ 0}
                                       +\colvec{2 \\ 1 \\ 0}\cdot y
                                       +\colvec{-1 \\ 0 \\ 1}\cdot z
                             \suchthat y,z\in\Re}
          \end{equation*}
        \item 参数化给出 $x=-(1/2)-(1/2)y-2z$，
          所以这就是向量描述。
          \begin{equation*}
            \colvec{x \\ y \\ z}=\colvec{-1/2 \\ 0 \\ 0}
                                       +\colvec{-1/2 \\ 1 \\ 0}\cdot y
                                       +\colvec{-2 \\ 0 \\ 1}\cdot z
          \end{equation*}
        \item 这里 
          $x=10-y-z-w$，于是我们得到含三个参数的向量描述。
          \begin{equation*}
            \colvec{x \\ y \\ z \\ w}=\colvec{10 \\ 0 \\ 0 \\ 0}
                                       +\colvec{-1 \\ 1 \\ 0 \\ 0}\cdot y
                                       +\colvec{-1 \\ 0 \\ 1 \\ 0}\cdot z
                                       +\colvec{-1 \\ 0 \\ 0 \\ 1}\cdot w
          \end{equation*}
      \end{exparts}
    \end{answer}
  \item 
    描述包含这个点与这条直线的平面。
    \begin{equation*}
      \colvec[r]{2 \\ 0 \\ 3}
      \qquad
      \set{\colvec[r]{-1 \\ 0 \\ -4}
           +\colvec[r]{1 \\ 1 \\ 2}t
           \suchthat t\in\Re}
    \end{equation*}
    \begin{answer}
      向量
      \begin{equation*}
        \colvec[r]{2 \\ 0 \\ 3}
      \end{equation*}
      不在这条直线上。
      由于
      \begin{equation*}
        \colvec[r]{2 \\ 0 \\ 3}
        -\colvec[r]{-1 \\ 0 \\ -4}
        =\colvec[r]{3 \\ 0 \\ 7}
      \end{equation*}
      我们可以这样描述该平面。
      \begin{equation*}
        \set{\colvec[r]{-1 \\ 0 \\ -4}
             +m\colvec[r]{1 \\ 1 \\ 2}
             +n\colvec[r]{3 \\ 0 \\ 7}
            \suchthat m,n\in\Re}
      \end{equation*}  
   \end{answer}
  \recommended \item
    求这两个平面的交。
    \begin{equation*}
      \set{\colvec[r]{1 \\ 1 \\ 1}t+
           \colvec[r]{0 \\ 1 \\ 3}s
           \suchthat t,s\in\Re}
      \qquad
      \set{\colvec[r]{1 \\ 1 \\ 0}
           +\colvec[r]{0 \\ 3 \\ 0}k+
           \colvec[r]{2 \\ 0 \\ 4}m
           \suchthat k,m\in\Re}
    \end{equation*}
    \begin{answer}
      重合的点就是下面这个方程组的解。
      \begin{equation*}
        \begin{linsys}{2}
         t  &  &   &= &1+2m\hfill      \\
         t  &+ &s  &= &1+3k\hfill      \\
         t  &+ &3s &= &4m\hfill
        \end{linsys}
      \end{equation*}
      用高斯消元法
      \begin{align*}
        \begin{amat}{4}
          1  &0  &0  &-2  &1  \\
          1  &1  &-3 &0   &1  \\
          1  &3  &0  &-4  &0
        \end{amat}
        &\grstep[-\rho_1+\rho_3]{-\rho_1+\rho_2}
        \begin{amat}{4}
          1  &0  &0  &-2  &1  \\
          0  &1  &-3 &2   &0  \\
          0  &3  &0  &-2  &-1
        \end{amat}                       \\                         
        &\grstep{-3\rho_2+\rho_3}
        \begin{amat}{4}
          1  &0  &0  &-2  &1  \\
          0  &1  &-3 &2   &0  \\
          0  &0  &9  &-8  &-1
        \end{amat}
      \end{align*}
      得到 \( k=-(1/9)+(8/9)m \)，所以 \( s=-(1/3)+(2/3)m \) 且 \( t=1+2m \)。
      交是下面这个集合。
      \begin{equation*}
        \set{\colvec[r]{1 \\ 1 \\ 0}+
             \colvec[r]{0 \\ 3 \\ 0}(-\frac{1}{9}+\frac{8}{9}m)+
             \colvec[r]{2 \\ 0 \\ 4}m
             \suchthat m\in\Re}
        =\set{\colvec[r]{1 \\ 2/3 \\ 0}
             +\colvec[r]{2 \\ 8/3 \\ 4}m
             \suchthat m\in\Re}
      \end{equation*}    
    \end{answer}
  \recommended \item
    在可能的情况下，求每一对的交。
    \begin{exparts}
      \partsitem \( \set{\colvec[r]{1 \\ 1 \\ 2}+t\colvec[r]{0 \\ 1 \\ 1}
                     \suchthat t\in\Re} \),\,
            \( \set{\colvec[r]{1 \\ 3 \\ -2}+s\colvec[r]{0 \\ 1 \\ 2}
                     \suchthat s\in\Re} \)
      \partsitem \( \set{\colvec[r]{2 \\ 0 \\ 1}+t\colvec[r]{1 \\ 1 \\ -1}
                     \suchthat t\in\Re} \),\,
            \( \set{s\colvec[r]{0 \\ 1 \\ 2}
                     +w\colvec[r]{0 \\ 4 \\ 1}
                     \suchthat s,w\in\Re} \)
    \end{exparts}
    \begin{answer}
       \begin{exparts}
        \partsitem 方程组
          \begin{equation*}
            \begin{linsys}{1}
              1     &= &1\hfill         \\
              1+t   &= &3+s\hfill         \\
              2+t   &= &-2+2s\hfill
            \end{linsys}
          \end{equation*}
          给出 \( s=6 \) 与 \( t=8 \)，所以这就是解集。
          \begin{equation*}
            \set{\colvec[r]{1 \\ 9 \\ 10}     }
          \end{equation*}
        \partsitem 这个方程组
          \begin{equation*}
            \begin{linsys}{1}
            2+t   &=  &0 \hfill        \\
            t     &=  &s+4w\hfill        \\
            1-t   &=  &2s+w\hfill
            \end{linsys}
          \end{equation*}
          给出 \( t=-2 \)、\( w=-1 \) 与 \( s=2 \)，所以它们的交
          是这个点。
          \begin{equation*}
            \colvec[r]{0 \\ -2 \\ 3}
          \end{equation*}
      \end{exparts}  
    \end{answer}
  \item 
      我们应当如何定义 $\Re^0$？
      \begin{answer}
        我们将在后面把它定义成恰有一个元素的集合\Dash 一个
        “原点”。  
      \end{answer}
  \puzzle \item   
    \cite{MathMag57p173}
    一个人以每小时
    \( 3 \) 英里的速率向东行进，发觉风好像从正北方向吹来。
    当他的速度加倍时，风好像来自东北方向。
    风的速度是多少？
    \begin{answer}
      \answerasgiven %
      向量三角形如下图所示，因此 \( \vec{w}=3\sqrt{2} \)，
      风从西北方向吹来。
      \begin{center}
        \setlength{\unitlength}{4pt}      % wind triangles.
        \begin{picture}(20,12)(0,0)
           \put(0,10){\vector(1,0){20} }
           \put(0,10){\vector(1,-1){10} }
              \put(3.8,5){\makebox(0,0)[r]{\small \( \vec{w} \)} }
           \put(10,0){\vector(1,1){10} }
           \put(10,0){\line(0,1){10} }
        \end{picture}
      \end{center}  
     \end{answer}
  \item \checked \label{exer:Euclid}  
    欧几里得（Euclid）把平面描述为
    “一个与其上的直线处处均匀相合的面”。
    希罗（Heron）等注释家把它解释为：
    “（平的面是）这样的面：若一条直线经过 
    其上的两点，
    则这条直线处处、以一切方式完全与该面重合”。
    （译文取自 \cite{Heath}，第171-172页。）
    本节中所描述的平面具有这种性质吗？
    这个描述足以恰当地定义平面吗？
    \begin{answer}
      欧几里得无疑是在设想 \( \Re^3 \) 之内的平面。
      然而，请注意 \( \Re^1 \) 与 \( \Re^2 \) 也都满足
      这个定义。  
    \end{answer}
\end{exercises}
















